% 10 - CONCLUSION
% ==============================================================
\sectionintro{11 · Synthèse}{Conclusion}{Bilan de l'architecture, de la PLC Visualization V1, du relevé d'essais et des limites identifiées.}

\twocolpage{
  \begin{primarycard}
    \textbf{Résultat obtenu}\par\small
    Le carrousel est structuré autour de trois modes et d'une machine d'états dédiée au cycle \texttt{AUTO}. La logique mémorise les huit compteurs, réalise quatre cycles de commande avant/arrière par indexation et passe en \texttt{ATTENTE\_BAC} lorsque le bac suivant est déjà plein. La PLC Visualization V1 expose le suivi et les commandes de simulation.
  \end{primarycard}

  \subsection{Résultats des essais}
  \begin{itemize}\small
    \item passage de \texttt{INIT} à \texttt{AUTO} après validation du zéro ;
    \item comptage sur front montant du capteur laser ;
    \item quatre cycles temporisés et bouclage logique du bac 8 vers le bac 1 ;
    \item réinitialisation individuelle ou globale des compteurs ;
    \item traitement de trois situations limites ;
    \item acquittement du défaut avec retour par \texttt{INIT}.
  \end{itemize}
}{
  \subsection{Démarche de validation}
  \small
  Le relevé existant couvre le fonctionnement nominal, la saturation des huit bacs et deux paramètres invalides. La synchronisation documentaire a ensuite comparé ces affirmations à \texttt{MAIN.TcPOU}, \texttt{FB\_Carrousel.TcPOU}, aux DUTs, à la Visualization, à son gestionnaire, à la liste de textes et au projet \texttt{.plcproj}.

  \begin{card}
    \textbf{Version livrée}\par\small
    Les sources de référence sont les objets TwinCAT du dépôt : \texttt{POUs/MAIN.TcPOU}, \texttt{POUs/FB\_Carrousel.TcPOU}, les deux DUTs et les objets de PLC Visualization. Aucun fichier agrégé \texttt{FB\_Carrousel.st} n'est livré.
  \end{card}

  \begin{notebox}[colback=eksuccess,colframe=green!30]{Synthèse}
    Le rapport décrit désormais l'implémentation et la Visualization V1 sans leur attribuer les fonctions absentes : retours de fin de course, sécurité machine, alarmes avancées, historique ou TwinCAT HMI V2.
  \end{notebox}
}

\vfill
{\color{ekline}\rule{\linewidth}{0.6pt}}\vspace{3mm}
\eyebrow{Références techniques}
\begin{itemize}\footnotesize\color{ekmuted}
  \item Beckhoff Information System, \href{https://infosys.beckhoff.com/content/1033/tcplclib_tc2_standard/74391563.html}{\texttt{R\_TRIG} - détection d'un front montant dans la bibliothèque Tc2\_Standard}.
  \item Beckhoff Information System, \href{https://infosys.beckhoff.com/content/1033/tc3_plc_intro/7647315979.html}{exemple TwinCAT 3 associant machine d'états, temporisateur et trigger}.
\end{itemize}

\newpage

% ==============================================================
